\documentclass[../main.tex]{subfiles}
\graphicspath{{\subfix{../images/}}}
\begin{document}

\subsection{Summary of Contributions}

We present three key advances: (1) a sliding window framework addressing boundary limitations with configurable overlap, (2) successful real-world validation on PLY files (61,914 points), and (3) comprehensive comparative analysis establishing data-dependent performance. The main finding: sliding window excels on noisy real data despite lower absolute magnitudes, providing 3.5× higher resolution and superior visual contrast through enhanced background suppression.

\subsection{Practical Impact}

\subsubsection{Method Selection Guidelines}

We established evidence-based guidelines for method selection.
\begin{itemize}
    \item \textbf{Real-world applications:} Sliding window recommended for superior noise handling
    \item \textbf{High-speed requirements:} Fixed patches for computational efficiency
    \item \textbf{Precision localization:} Sliding window for detailed spatial analysis
\end{itemize}

\subsubsection{Theoretical Justification}

Our work validates Shape Subspace methods for anomaly detection: subspaces capture complete geometric patterns (not just height variations), SVD naturally filters noise while preserving features, and Difference Subspace provides a mathematically rigorous comparison framework more interpretable than deep learning and more sophisticated than simple statistical methods.

\subsection{Future Work}

\subsubsection{Immediate Extensions}

\textbf{Adaptive Parameter Selection:} Develop algorithms for automatic selection of the percentage of overlap and the size of the window based on the density of local points and noise characteristics.

\textbf{Multi-Scale Analysis:} Implement hierarchical sliding window analysis with varying window sizes to capture anomalies at different scales simultaneously.

\textbf{Performance Optimization:} Investigate parallel processing strategies and GPU acceleration for real-time analysis of large point clouds.

\subsubsection{Practical Applications}

\textbf{Infrastructure Monitoring:} Apply the method to the monitoring of structural health health of bridges, roads, and buildings for automated defect detection.

\textbf{Quality Control:} Implement in manufacturing environments for inspection of surface quality of components produced.

\textbf{Autonomous Systems:} Integrate with autonomous vehicle perception systems for real-time road condition assessment.

\subsection{Conclusion}

This progress report demonstrates substantial advancement from synthetic validation to real-world application. The sliding window extension addresses boundary limitations while maintaining mathematical rigor. Our key finding, that method effectiveness is data dependent, provides practical guidance: use fixed patches for clean data and sliding windows for noisy real-world data. This successful transition to real PLY files represents a crucial step toward practical deployment of shape subspace methods for surface anomaly detection.

\end{document}